\section{Related Work}\label{sec:related_work}
Software quality and maintainability have been extensively studied by the research community, covering a wide range of topics such as code quality metrics~\cite{mccabe_complexity_1976, chidamber_metrics_1994, mo_decoupling_2016}, change-proneness~\cite{chowdhury_revisiting_2022, khomh_exploratory_2012, romano_using_2011}, bug-proneness~\cite{chowdhury_method-level_2024, pascarella_performance_2020, giger_method-level_2012}, code readability~\cite{posnett_simpler_2011, buse_learning_2010}, and code understandability~\cite{scalabrino_automatically_2017}, among others. Due to its relevance to software quality and maintainability, a different stream of research has focused on SATD detection, classification, and its long-term maintenance impact. 

Potdar and Shihab~\cite{potdar_exploratory_2014} were the first to systematically study SATD in four large open-source projects, without differentiating between production and test code. They found SATD in 2.4–31.0\% of project files and investigated the reasons for its introduction and its removal frequency. Bavota and Russo~\cite{bavota_large-scale_2016} replicated the study of Potdar and Shihab, but by analyzing a large number of projects (159). They found that SATD comments tend to increase over time and persist for long periods. Chowdhury et al.~\cite{chowdhury_evidence_2025} conducted a large-scale method-level analysis of 774,051 methods across 49 open-source projects, revealing that methods with SATD are larger, more complex, less readable, and more prone to bugs and frequent changes. Similarly, Wehaibi et al.~\cite{wehaibi_examining_2016} analyzed five major open-source systems and found that technical debt negatively affects software evolvability, making future changes more difficult. Collectively, these studies underscore the negative impact of SATD on long-term software health, highlighting the importance of its classification and detection.

The detection of SATD has been explored extensively through both manual and automated approaches. Potdar and Shihab~\cite{potdar_exploratory_2014} identified 62 frequent keywords showing their potential in SATD detection. Various automated techniques have been proposed as well. Guo et al.~\cite{guo_how_2021} introduced the Matches task Annotation Tags (MAT) method, which achieves performance comparable to NLP-based approaches~\cite{maldonado_using_2017}. Huang et al.~\cite{huang_identifying_2018} applied text mining to extract features, performed feature selection, and trained a composite classifier for SATD detection. More recently, large language models (LLMs)~\cite{sheikhaei_empirical_2024, lambert_identification_2024,gu_self-admitted_2024} have been employed, outperforming traditional methods for detecting SATD in application code.

To better conceptualize technical debt, Alves et al.~\cite{alves_towards_2014} proposed an ontology synthesizing definitions and indicators scattered across the literature, categorizing technical debt into 13 types (e.g., architecture, build, code, test). Building on this foundation, Maldonado and Shihab~\cite{maldonado_detecting_2015} manually reviewed 33,093 heuristically filtered source code comments from five Java projects and identified 2,457 SATD comments, which they classified into five categories: design, requirement, defect, test, and documentation. Maldonado et al.~\cite{maldonado_using_2017} later extended this study to additional projects and applied natural language processing (NLP) techniques for classification. Their dataset has been used by subsequent studies on automated SATD detection and classification \cite{ sheikhaei_empirical_2024, guo_how_2021, gu_self-admitted_2024, arcelli_fontana_binary_2025, fucci_waiting_2021, njeru_hybrid_2025}.
Recent work has expanded SATD research into new domains beyond traditional software projects. OBrien et al.~\cite{obrien_artifact_2022} categorized SATD comments in machine learning code into 23 distinct types, while Bhatia et al.~\cite{bhatia_empirical_2025} extended Bavota et al.’s~\cite{bavota_large-scale_2016} taxonomy by adding new machine-learning-specific categories and subcategories to classify SATD in machine learning software. Liu et al.~\cite{liu_is_2020} examined seven popular open-source deep learning frameworks to assess SATD prevalence and categories. Pinna et al.~\cite{pinna_investigation_2023} found that SATD in blockchain projects is more prevalent than in non-blockchain systems. Azuma et al.~\cite{azuma_empirical_2022} manually classified SATD comments in Dockerfiles, identifying five main classes—including Docker-specific categories—to capture domain-specific debt. Similarly, Wilder et al.~\cite{wilder_exploratory_2023} analyzed 15,614 open-source Android applications and found that code debt is the most prevalent type.

Although a study by Counsell and Swift~\cite{counsell_timing_2022} examined SATD in test code in addition to source code, their focus was on understanding whether the time of day is related to writing SATD, rather than on identifying patterns of test code SATD or supporting their detection. During the course of our study, we became aware of a parallel work by Nakamura et al.~\cite{nakamura2025understanding}, which also investigated SATD in test code. However, our study differs substantially from theirs in several important aspects, as outlined below.

\begin{itemize}
     \item To identify test-code SATD, we manually examined 50,000 code comments. In contrast, the other study relied on an existing SATD detection tool, DebHunter, which was not specifically designed or trained to detect SATD in test code. Our evaluation of DebHunter on test-code SATD detection revealed low accuracy (RQ2). As such, using DebHunter to construct ground-truth datasets, as done in ~\cite{nakamura2025understanding}, may have introduced inaccuracies into the identified SATD instances and affected the reliability of the resulting analysis.    
     \item While constructing their taxonomy of test-code SATD, Nakamura et al. did not attempt to align their findings with existing SATD taxonomies. Such comparisons are important because they help determine how newly identified categories differ from or extend established ones, and this practice is commonly followed in software engineering research. For example, in their work, ``Why Attention Fails: A Taxonomy of Faults in Attention-Based Neural Networks'', Jahan et al.~\cite{jahan2025attention} introduced new categories only when the identified issues could not be adequately mapped to an existing deep-learning bug taxonomy. Similarly, in our study, we first attempted to map each SATD instance to the existing source-code SATD taxonomy and introduced new categories only when the comments could not be appropriately classified within the existing structure. Without such an analysis, it remains unclear whether a separate taxonomy for test-code SATD is necessary or whether existing SATD taxonomies are already sufficient to capture these instances.
    \item We argue that the problem of identifying whether a comment constitutes SATD should be addressed before attempting to classify it into specific SATD categories. If SATD comments cannot be reliably distinguished from non-SATD comments, then classifying their types becomes considerably less meaningful. However, the competing study focused on the multiclass classification task (i.e., predicting SATD categories) without first addressing the more fundamental binary classification task of distinguishing SATD from non-SATD comments. This concern is particularly significant because our results show that existing SATD detection tools, as well as both open-source and proprietary LLMs, struggle to accurately identify SATD in test code. These findings suggest that reliable test code SATD detection remains an open challenge and should be established before pursuing prediction of SATD classes.
    \item In addition to evaluating a wide range of open-source and proprietary LLMs, we also evaluated state-of-the-art SATD detection tools. This comparison is motivated by the practical consideration that LLM-based solutions can be computationally expensive, and therefore should ideally be used only if simpler approaches do not provide sufficient performance.


\end{itemize}

\begin{tcolorbox}
\textbf{\underline{Motivation}:}
While numerous studies have examined different aspects of test code maintenance \cite{spadini_relation_2018, lin_quality_2019, beller_developer_2019, beller_when_2015, zaidman_studying_2011, athanasiou_test_2014}, investigating the detection and classification of SATD in test code remains underexplored. Motivated by this gap, we conduct an empirical study on the test code SATD, exploring its types and detection.
\end{tcolorbox}

% . Spadini et al.~\cite{spadini_relation_2018} found that test code containing smells is more prone to changes and defects. Lin et al.~\cite{lin_quality_2019} analyzed manually written and automatically generated test cases from ten open-source projects and observed that test code often contains low-quality identifiers—even in manually written tests—and that identifier quality is generally lower than in production code.

% The most relevant works to this study are those by Potdar and Shihab \cite{potdar_exploratory_2014}, who identified SATD comments by manually reviewing source code comments, and by Maldonado and Shihab \cite{maldonado_detecting_2015}, who conducted a manual analysis of comments from five open-source Java projects and identified five distinct categories of technical debt articulated. Similarly, Bavota and Russo \cite{bavota_large-scale_2016} manually categorized 366 comments from 159 software projects into different categories and sub-categories and Wilder II et al. \cite{ii_exploratory_2023} categorized 386 randomly selected SATD comments from Android applications.


% TODO : add a paragraph showing that SATd negatively impact software maintenace
% explain TODO: Similar to LLM paper
%  Negative Effect of SATD

% Removal of SATD

%TODO : add proof: achieving superior performance compared to traditional rule-based techniques.
